\selectlanguage{slovak}

\subsection*{Úvod}
Táto diplomová práca skúma uskutočniteľnosť vybudovania praktického, offline asistenta pre príkazový riadok (\gls{cli}) systému Linux (shAI) s~využitím malých lokálnych modelov \gls{llm} na bežne dostupnom spotrebiteľskom hardvéri. Zameriava sa na~súkromie, dostupnosť bez pripojenia k~sieti a~interaktívnu latenciu v~podmienkach obmedzených zdrojov.
Výsledný prototyp, shAI, je open-source a~slúži ako referenčný základ pre budúci výskum. Nižšie uvedený prehľad stručne zhŕňa použité metódy, implementáciu a~empirické výsledky.

% Práca si stanovuje tri hlavné ciele:
% \begin{enumerate}
%   \item Analyzovať prostredie existujúcich terminálových asistentov a~identifikovať ich obmedzenia pre lokálne použitie.
%   \item Navrhnúť a~implementovať prototyp systému, ktorý funguje výhradne na lokálnych modeloch.
%   \item Vyhodnotiť uskutočniteľnosť kľúčových komponentov pomocou objektívnych metrík.
% \end{enumerate}


\subsection*{Teoretické zázemie}

Táto sekcia popisuje teoretické zázemie relevantné pre zvolenú tému.

\subsubsection*{Koncepty terminálových rozhraní}

Táto sekcia definuje rozhrania relevantné pre naše riešenie: \gls{cli} a~\gls{tui}.

\paragraph{Rozhranie} je zdieľaná hranica, na ktorej si dva samostatné komponenty vymieňajú informácie určitým definovaným spôsobom~\cite{hookway2014}.

\paragraph{Hardvérový terminál} je elektronické zariadenie na zadávanie a~prepisovanie dát, ktoré poskytuje textovú interakciu s~počítačovými systémami~\cite{terminfo}.

\paragraph{Emulátor terminálu} (ďalej len \textbf{terminál}) replikuje funkčnosť hardvérového terminálu vo forme softvéru.

\paragraph{Terminálová aplikácia} beží v~rámci emulátorov terminálu a~poskytuje textové používateľské rozhrania prostredníctvom \gls{cli} alebo \gls{tui}. Komunikácia prebieha cez štandardné vstupné a~výstupné prúdy, čo obmedzuje interakciu iba na~text~\cite{man_man,linfo_terminal}.

\paragraph{\gls{cli}} rozširuje definície rozhraní na rozhrania príkazového riadka –~ide o~textové rozhrania, kde používatelia zadávajú príkazy a~dostávajú textové odpovede.

\paragraph{Shell} je program poskytujúci priamy prístup k~systému, ktorý implementuje \gls{cli}, \gls{tui} alebo \gls{gui}. \gls{cli} shelly umožňujú používateľom spúšťať príkazy a~\gls{cli} nástroje~\cite{linfo_terminal}.

\paragraph{\gls{cli} nástroj} (angl. \textsl{utility}) je program spúšťaný v~prostrediach \gls{cli}.

\paragraph{\gls{tui}} rozširuje \gls{cli} tým, že poskytuje štruktúrované, intuitívne rozhrania s~využitím knižníc ako ncurses a~rozšírených funkcionalít terminálu (angl. \textsl{terminal capabilities}) ~\cite{linfo_terminal}.
\subsection*{Koncepty \gls{llm}}

\paragraph{Jazykové modely} predpovedajú a~generujú textové sekvencie, čo je kľúčové pre preklad, sumarizáciu a~tvorbu obsahu. Ich vývoj napreduje spolu s~pokrokmi v~oblasti neurónových sietí~\cite{naveed2023}.

\paragraph{\gls{llm}} sú jazykové modely založené na architektúre transformer s~veľkým počtom parametrov, trénované na veľkých textových datasetoch (napr. \gls{gpt}~\cite{radford2018improving}, \gls{bert}~\cite{devlin2018bert}, \gls{lstm}~\cite{hochreiter1997lstm} alebo Mamba~\cite{gu2023mamba}). Vykazujú dobré porozumenie prirodzenému jazyku a~schopnosť riešiť komplexné úlohy prostredníctvom generovania textu.

\subsubsection*{Architektúra \gls{llm}}

Architektúru \gls{llm} možno opísať vo vrstvách, pričom každá z~nich spracováva vstupnú informáciu a~odovzdáva ju nasledujúcej vrstve~\cite{mikolov2013distributed,pennington2014glove,zhao2023, attention}:

\begin{itemize}
    \item \textbf{Tokenizátor} rozkladá vstupný text na tokeny s~fixnou slovnou zásobou definovanou počas trénovania.
    \textsl{Nové alebo zriedkavé slová nemusia byť spracované presne v~závislosti od veľkosti tokenizátora.}

    \item \textbf{Embeddingová vrstva} prevádza tokenizovaný vstup na vektorové reprezentácie (ako sú fixné vektory alebo vektory závislé od pozície).

    \item \textbf{Transformerové vrstvy} pozostávajú z~viacerých \textsl{attention heads}~\cite{attention} a~dopredných (feedforward) sietí, ktoré modelu umožňujú zachytiť komplexné vzťahy medzi tokenmi.

    \item \textbf{Výstupná vrstva} generuje finálny výstup (text, pravdepodobnosti) prostredníctvom lineárnej vrstvy s~aktiváciou softmax.
\end{itemize}

V~pôvodnom článku sú na báze mechanizmu pozornosti postavené dva mechanizmy~\cite{zhao2023}:
\begin{itemize}
    \item \textbf{Kodér} mapuje vstupné sekvencie tokenov na spojité reprezentácie pomocou \textsl{multi-head attention} a~dopredných sietí.
  \item \textbf{Dekodér} autoregresívne generuje výstupné sekvencie tokenov. Zahŕňa subvrstvy v~štýle kodéra a~navyše \textsl{multi-head attention} nad výstupom kodéra.
\end{itemize}

Moderné modely \gls{llm} využívajú primárne architektúru typu len s dekodérom (angl. \textsl{decoder-only architecture}), ktorá sa spravidla ľahšie škáluje, trénuje a~dotrénuje (angl. \textsl{fine-tune}) než modely typu Kodér-Dekodér (angl. \textsl{encoder-decoder}), pričom dosahuje lepší výkon pri rovnakom počte trénovacích tokenov~\cite{zhao2023,wei2023chainofthoughtpromptingelicitsreasoning}.

\paragraph{Kontextové okno} je maximálny počet tokenov, ktoré dokáže Transformer spracovať naraz. Kratšie prompty sa spracovávajú rýchlejšie; príliš dlhé prompty strácajú informácie, ak prekročia kapacitu okna.

\paragraph{\gls{moe}}
Architektúra rozdeľuje doprednú sieť do viacerých špecializovaných subneurónových sietí (expertov), pričom smerovač (router) vyberá na spracovanie iba ich časť (zvyčajne jedného alebo dvoch).
Modely typu \gls{moe} sú výpočtovo efektívnejšie, ale dosahujú horší výkon ako rovnako veľké husté modely s~rovnakým celkovým počtom parametrov. Avšak v~prípadoch, keď nie je obmedzená pamäť VRAM, môže \gls{moe} prekonať výkon hustých modelov pri zachovaní ekvivalentných výpočtových nákladov~\cite{dai2024deepseekmoe}.

\paragraph{Embedding \gls{llm}} sú neurónové modely založené na architektúre transformer, ktoré slúžia na~generovanie vektorových reprezentácií dokumentov.
Zachytávajú komplexnejšie vzťahy medzi tokenmi než tradičné modely, pričom využívajú mechanizmy pozornosti na váženie dôležitosti tokenov, vďaka čomu poskytujú presnejšie sémantické reprezentácie~\cite{zhang2024systematicsurveytextsummarization,khan2026retrieval}.

\subsection*{Pokročilé koncepty \gls{llm}}

Táto sekcia spomína dôležité koncepty \gls{llm}, ktoré sa bežne využívajú v~rôznych aplikáciách postavených na \gls{llm}.

\subsubsection*{Systémy AI agentov založené na \gls{llm}}

\paragraph{AI agenti} sú autonómne entity schopné plánovania, rozhodovania a~vykonávania akcií s~cieľom dosiahnuť komplexné ciele~\cite{naveed2023,zhao2023}.

AI agenti prinášajú také vylepšenia pre \gls{llm}:
\begin{itemize}
  \item \textbf{Hierarchické plánovanie:} Rozkladajú zložité úlohy na zvládnuteľné čiastkové úlohy.
  \item \textbf{Používanie nástrojov:} Využívajú externé nástroje, RAG pipeline alebo API nad rámec samotného generovania textu.
  \item \textbf{Iteratívne generovanie:} Generujú viacero krokov uvažovania iteratívne, pričom zjemňujú výstupy a~prispôsobujú sa novým informáciám.
\end{itemize}
Pracovný postup typického \gls{llm} agenta popisuje Algoritmus~\ref{alg:agent}.

AI agenti dekomponujú komplexné ciele do zvládnuteľných krokov, pričom iteratívne dopytujú samých seba s~využitím rôznych nástrojov~\cite{nvidia-llm-agent}.

\subsubsection*{Retrieval-Augmented Generation (RAG)}

\textbf{\gls{rag}} (vyhľadávaním rozšírené generovanie) je technika, ktorá umožňuje modelom \gls{llm} využívať externé znalostné bázy, čím výrazne znižuje halucinácie a~zvyšuje presnosť~\cite{zhao2023}. Pracovný postup \gls{rag} popisuje Algoritmus~\ref{alg:rag}.

\paragraph{Granularita vyhľadávania}
Systémy \gls{rag} môžu vyhľadávať informácie na viacerých úrovniach granularity vrátane vyhľadávania \textbf{na úrovni dokumentov}, \textbf{na úrovni pasáží} a~\textbf{na úrovni viet}. Volba granularity ovplyvňuje výpočtovú záťaž pri vytváraní indexov a~vykonávaní vyhľadávacích operácií, presnosť vyhľadávania a~určuje veľkosť vygenerovaného výstupu.

Sekcia~\ref{sec:rag-eval} taktiež spomína rôzne metriky hodnotenia \gls{rag}, ako napríklad precision@k alebo MRR.

\subsubsection*{Fine-tuning}

Fine-tuning je proces, pri ktorom sa vezme predtrénovaný \gls{llm} a~prispôsobí sa konkrétnej úlohe dotrénovaním na menšom datasete špecifickom pre danú úlohu. Tento prístup umožňuje modelu naučiť sa špecifické vzorce a~zlepšiť svoj výkon na cieľovej úlohe pri zachovaní nízkeho úsilia a~nákladov na tréning~\cite{zhao2023}.

\subsection*{Reprezentácia a~vyhľadávanie dokumentov}

Táto sekcia poskytuje teoretické zázemie v~oblasti reprezentácie, ukladania a~vyhľadávania dokumentov.

\paragraph{Sémantická reprezentácia dokumentov} je proces konverzie dokumentov do formátu, ktorý zachytáva ich konceptuálny a~sémantický význam vo forme vhodnej pre strojové spracovanie. Sémantická reprezentácia dokumentov sa všeobecne delí do dvoch primárnych skupín~\cite{khan2026retrieval}:
\begin{itemize}
    \item \textbf{Riedké vektorové reprezentácie} využívajú tradičné prístupy založené na kľúčových slovách, ako napríklad \emph{bag-of-words} (\gls{bow}), one-hot kódovanie, latentné sémantické indexovanie (\gls{lsindexing}) alebo \emph{term frequency-inverse document frequency} (\gls{tfidf}). Tieto prístupy sú výpočtovo efektívne a~vhodné pre jednoduché dopyty, avšak môžu mať problémy pri komplexných dopytoch, ktoré si vyžadujú porozumenie kontextu~\cite{khan2026retrieval}.

    \item \textbf{Husté vektorové reprezentácie} využívajú embedding \gls{llm} na generovanie embedding vektorov, ktoré zachytávajú sémantický význam textu. Tento prístup je efektívnejší pre zložité dopyty, vyžaduje si však viac výpočtových zdrojov než riedke prístupy~\cite{khan2026retrieval}.

    \item Okrem toho sa \textbf{hybridné techniky vyhľadávania} pokúšajú kombinovať výhody oboch prístupov pomocou techník, ako je napríklad model SPLADE, ktorý obohacuje riedke reprezentácie o~sémanticky príbuzné výrazy~\cite{khan2026retrieval}.
\end{itemize}

Pre tento výskum sa ako primárny prístup zvolilo vyhľadávanie pomocou hustých reprezentácií, keďže je vhodnejšie pre komplexné dopyty a~dokáže zachytiť sémantický význam dokumentov, čo je kľúčové pre úlohu výberu vhodných dokumentov na základe používateľských dopytov.

\paragraph{Vektorové databázy} sú špecializované databázové systémy navrhnuté na ukladanie a~správu vysokorozmerných vektorových dát a~na vykonávanie efektívneho vyhľadávania podobnosti na základe rôznych metrík vzdialenosti vrátane kosínusovej podobnosti, L2 vzdialenosti a~euklidovskej vzdialenosti. Príkladmi široko používaných vektorových databáz sú FAISS (Facebook AI Similarity Search)~\cite{faiss2026}, Annoy~\cite{annoy2026} a~Milvus~\cite{milvus2026}. Bežne sa využívajú na ukladanie a~vyhľadávanie vektorových reprezentácií dokumentov.

\subsection*{Manuálové stránky systému Linux}

\textbf{Manuálové stránky systému Linux} (alebo \textbf{man pages}) sa spracovávajú ako dáta pre validáciu. \texttt{man} je systémový prehliadač manuálových stránok, ktorý zobrazuje dokumentáciu k~programom, nástrojom alebo funkciám~\cite{man_man}.

Manuálové stránky sa delia do skupín a~sekcií. Kľúčovými sekciami sú \textbf{DESCRIPTION} (podrobný popis programu, funkčnosti a~použitia) a~\textbf{OPTIONS} (prepínače, argumenty a~ich popisy)~\cite{man_man}.

\subsection*{Návrh a~implementácia}
Návrh vyžaduje fungovanie bez pripojenia k~sieti (offline) a~minimálne externé závislosti. Ako primárny zdroj znalostí boli zvolené manuálové stránky systému Linux.

Systém shAI je implementovaný ako modulárna aplikácia s~tromi komponentmi:
\begin{enumerate}
    \item Klasifikátor príkazov a~dopytov (nepouzitý vo finálnom riešení)
    \item Vyhľadávač dokumentov, ktorý vytvára a~prehľadáva indexy, založený na knižnici \gls{faiss} a~nástroji Ollama pre lokálne spúšťanie modelov
    \item Sumarizátor, ktorý generuje stručné odpovede z~vyhľadaného obsahu
\end{enumerate}
Rozhodnutia o~návrhu sú podporené diagramom architektúry a~diagramom prípadov použitia (pozri Obrázky~\ref{fig:use-cases} a~\ref{fig:component-arch}).

Okrem kľúčových komponentov zahŕňa implementácia aj používateľské procesy pre inicializáciu systému, zostavenie indexu, vykonávanie dopytov a~zmeny konfigurácie.

Komponent klasifikátora je navrhnutý ako prototyp. Môže byť vyvolaný nad aktuálnym používateľským vstupným bufferom a~vrátiť binárne rozhodnutie, avšak integrácia tohto riešenia do všeobecného shellu sa ukázala ako nepraktická. Nie je však bez potenciálu, keďže niektoré z~existujúcich riešení (napríklad \texttt{llm-zsh-plugin}) ukázali, že existujú špecifické spôsoby aplikácie, ako to dosiahnuť.

Komponent vyhľadávača (retriever) zabezpečuje zber dokumentov, generovanie embedding vektorov a~ukladanie vektorových indexov, pričom podporuje tri prístupy k~vyhľadávaniu:
\begin{enumerate}
  \item Vyhľadávanie na úrovni celých dokumentov
  \item Vyhľadávanie na úrovni dokumentov s~preusporiadaním (rerankingom) pasáží v~reálnom čase
  \item Vyhľadávanie na úrovni pasáží
\end{enumerate}
Umožňuje tiež explicitnú konfiguráciu rôznych parametrov vyhľadávania, ako sú top-k výsledky, veľkosť pasáže a~umiestnenie úložiska indexu, čo umožnilo vykonávanie riadených experimentov. Indexovanie sa vykonáva nad~vopred spracovanými manuálovými stránkami uloženými ako textové súbory v~lokálnom databázovom adresári.

Komponent sumarizátora poskytuje finálnu odpoveď používateľovi. Prijíma dopyt a~sadu vyhľadaných pasáží, naformátuje prompt s~explicitnou štruktúrou a~vygeneruje stručnú odpoveď zameranú na používateľskú otázku. Hoci bol pôvodne zamýšľaný ako plne funkčný agentový systém, počas počiatočných fáz manuálneho testovania sa ukázal ako nespoľahlivý a~nepresný, keďže modely \gls{llm} zvyčajne zlyhávali v~rozhodovacom procese, a~preto sa hlavný dôraz presunul na komponent vyhľadávača.

Riešenie využíva programovací jazyk Python s~klientskou knižnicou kompatibilnou s~OpenAI na prístup k~modelom, knižnicu \gls{faiss} na vyhľadávanie podobnosti a~štandardné systémové knižnice na správu súborového systému a~subprocesov. Lokálne runtime prostredie pre modely s~konfigurovateľnou dĺžkou kontextu a~podporou GPU zabezpečuje platforma Ollama. Testovacia sada je riadená konfiguračnými súbormi, čo umožňuje reprodukovateľné spúšťanie naprieč variantmi modelov a~stratégiami vyhľadávania.

Implementácia zahŕňa aj \gls{cli} vstupný bod, ktorý sprístupňuje načítanie konfigurácie a~spúšťanie testov. Konfiguračné súbory špecifikujú, ktorý model sa má použiť na embedding a~generovanie. Systém je navrhnutý tak, aby v~prípade potreby opätovne využil rovnaké rozhrania pre lokálne aj vzdialené backendy. Táto modulárna štruktúra umožňuje nezávislé testovanie komponentov vyhľadávania a~sumarizácie.

Rodiny modelov použité v~experimentoch sú zhrnuté v~Tabuľkách~\ref{tab:llm-summary} a~\ref{tab:embedding-summary}.


\subsection*{Vyhodnotenie}

Táto sekcia popisuje proces vyhodnotenia navrhnutých riešení.

\subsubsection*{Klasifikátor prirodzeného jazyka vs. príkazov shellu}

Na testovanie klasifikátora bola použitá ručne pripravená kolekcia \texttt{classifier-100}, ktorá obsahuje 100 vstupov s~vyváženým pomerom dopytov v~prirodzenom jazyku a~príkazov shellu. Je rozdelená nasledovne:
\begin{itemize}
  \item 50 dopytov v~prirodzenom jazyku (napr. „How do I~list files in a~directory?“)
  \item 50 vstupov vo forme príkazov shellu (napr. „ls -la /home/user“)
\end{itemize}

\paragraph{Výsledky testovania}{}
Tabuľka~\ref{tab:eval-classifier-results} zhŕňa presnosť klasifikácie a~priemernú latenciu na jeden dopyt pre každý model.

\paragraph{Hlavné zistenia a~obmedzenia}
Vyhodnotenie a~testovanie tohto komponentu odhalilo, že zvolený prístup ku klasifikácii je nespoľahlivý, pričom dosahuje presnosť iba 50~-- 87~\% s~vysokými latenciami až do 2,9~s, čo nie je dostatočne spoľahlivé a~rýchle pre obalovú vrstvu (wrapper) shellu. Navyše sa samotná koncepcia generického obalového segmentu (shell wrapper) stáva problemom a~výrazne komplikuje údržbu systému. Tento komponent zostáva vo fáze prototypu a~ďalej sa v~ňom nepokračovalo.

\subsubsection*{Vyhľadávač dokumentov (Document Retriever)}

\paragraph{Prístup k~vyhodnoteniu}
Vyhľadávač je vyhodnocovaný na základe jeho schopnosti vrátiť relevantné dokumenty v~rámci top-k výsledkov pomocou metrík Precision@k a~MRR, s~využitím troch variantov automatizovaných testov:
\begin{enumerate}
  \item Vyhľadávanie na úrovni celých dokumentov (test1)
  \item Vyhľadávanie na úrovni dokumentov s~preusporiadaním na úrovni pasáží (test2) \textsl{(najprv sa vyhľadajú dokumenty, následne sa preusporiadajú pasáže v~nich a~vyhodnocuje sa úspešnosť zhody pasáží)}
  \item Vyhľadávanie s~rozsahom na úrovni pasáží (test3)
\end{enumerate}
Testovanie prebehlo na testovacích kolekciách \texttt{sample-10} (14 dokumentov, 250 dopytov) a~\texttt{sample-100} (101 dokumentov, 300 dopytov).

\paragraph{Výsledky testovania}{}
Tabuľky~\ref{tab:eval-retriever-sample10-doc}~-- \ref{tab:eval-retriever-sample100-passage} zhŕňajú vyhľadávanie na úrovni dokumentov (test1), preusporiadanie na úrovni pasáží (test2) a~vyhľadávanie na úrovni pasáží (test3). Všetky testy používali parameter top\_k=5; test2 využíva top\_k\_extended=10.

\paragraph{Hlavné zistenia}
Komponent vyhľadávača dokumentov sa ukázal ako robustný a~prakticky použiteľný pri všetkých prístupoch z~hľadiska presnosti. Na kolekcii \texttt{sample-10} sa presnosť Top-1 pohybovala v~rozmedzí 83,6~-- 86,8~\%, pričom presnosť Top-5 dosahovala až 100~\%. Na kolekcii \texttt{sample-100} sa Top-1 pohybovala od 62,7~\% do 72,7~\% a~Top-5 dosahovala až 92,0~\%. Model qwen3-embedding:0.6b poskytol najlepšiu presnosť, avšak čas zostavenia jeho indexu bol na väčšom korpuse podstatne vyšší (128~-- 134~s v~porovnaní s~1,6~-- 2,3~s pri embedding Granite). Je však dôležité poznamenať, že druhý a~tretí prístup vykazujú omnoho menšiu veľkosť výstupu (v priemere 5-tisíc znakov, čo je približne 4-tisíc tokenov, oproti priemerne 80-tisíc znakom pri prvom prístupe). To je kritické pre komponent sumarizátora, pretože to výrazne znižuje latenciu a~zvyšuje presnosť generovanej odpovede.

\subsubsection*{End-to-end testovanie}

\paragraph{Prístup k~vyhodnoteniu}
Kompletný systém \gls{rag} (vyhľadávanie dokumentov +~sumarizácia informácií) je vyhodnocovaný end-to-end s~cieľom posúdiť reálny výkon vrátane času indexovania dokumentov, spracovania dopytov a~latencie odpovede používateľovi pri rôznych hardvérových konfiguráciách a~veľkostiach korpusu dokumentov. Hodnotenie sa zameriava na~praktickú použiteľnosť pre pomoc v~\gls{cli}.

\paragraph{Testovacia kolekcia}{}
Pre end-to-end testovanie používame sadu dopytov (10 dopytov) a~spúšťame ju proti kolekciám \texttt{sample-10}, \texttt{sample-100} a~\texttt{sample-1000}.

\paragraph{Výsledky testovania}{}
Tabuľky~\ref{tab:eval-21-results}~-- \ref{tab:eval-32-results} predstavujú ukážkové výsledky z~vyhodnotenia pre každú kolekciu. Hodnoty reprezentujú priemerné latencie na dopyt v~sekundách. Vyhodnotenie sa uskutočnilo na systéme s~grafickou kartou NVIDIA RTX 4060 8GB. Testovanie bolo vykonané pre najľahšie (qwen2.5:0.5b pre generovanie a~ibm/granite-embedding:30m pre vyhľadávanie) a~najťažšie (qwen3:1.7b pre generovanie a~qwen3-embedding:0.6b pre vyhľadávanie) modely, aby sa poskytol rozsah očakávaného výkonu naprieč rôznymi hardvérovými konfiguráciami.

\paragraph{Hlavné zistenia}
Vyhľadávanie na úrovni dokumentov a~vyhľadávanie na úrovni dokumentov s~preusporiadaním pasáží vykazujú najlepší výkon pri použití malých embedding modelov, ako je ibm/granite-embedding:30m. Latencia preusporiadania pasáží v~reálnom čase je dostatočne nízka na~to, aby fungovala lepšie než prístup na úrovni pasáží, pričom zachováva malú veľkosť výstupu a~rovnakú presnosť. Naopak, pri použití väčších embedding modelov, ako je qwen3-embedding:0.6b, prístup na úrovni pasáží prekonáva ostatné prístupy, keďže výpočtový čas potrebný na tvorbu vektorov pre celý dokument alebo preusporiadanie v~reálnom čase sa stáva príliš nákladným.

\subsubsection*{Regresné testy}
Implementované softvérové regresné testy slúžia na~zaistenie stability a~spoľahlivosti systému shAI počas vývoja a~budúcej údržby. Testy pokrývajú kľúčové vetvy aplikácie vrátane vyhľadávania dokumentov, sumarizácie a~end-to-end spracovania dopytov. Spolu so softvérovými regresnými testami sme implementovali aj sadu regresných testov prostredia, ktorá zabezpečuje, že pri akejkoľvek zmene systému nedôjde k~narušeniu žiadneho z~testov (najmä štatistických testov, ktorých spustenie zvyčajne vyžaduje veľa času). To nám umožňuje s~ľahkosťou udržiavať stabilitu systému a~rýchlo identifikovať akékoľvek problémy, ktoré by mohli nastať počas vývoja.

\subsubsection*{Diskusia o~budúcich smeroch vývoja}

Budúca práca v~rámci samotnej aplikácie navrhuje tieto vylepšenia:
\begin{itemize}
    \item \textbf{Širšia podpora úloh:} Testovať iné ohraničené úlohy, ako napríklad vysvetľovanie chýb alebo generovanie na základe šablón.
    \item \textbf{Vylepšenia používateľského rozhrania:} Pridať bohatšiu interakciu pri zachovaní nízkej latencie.
    \item \textbf{Optimalizácia softvéru:} Zlepšiť využitie cache, načítavanie modelov a~optimalizovanie vyhľadávania.
    \item \textbf{Metodológia:} Porovnať metódy vyhodnocovania a~preskúmať destiláciu, kvantizáciu alebo novšie lokálne architektúry.
\end{itemize}
Budúci výskum nad rámec vyvinutej aplikácie môže pokračovať v~týchto smeroch:
\begin{itemize}
    \item \textbf{Porovnávanie rôznych prístupov a~frameworkov na testovanie modelov}
    \item \textbf{Optimalizácia modelov a~skúmanie architektúr}
    \item \textbf{Skúmanie rôznych prístupov k~riadeniu a~prepájaniu modelov (harnessing)}
\end{itemize}

\subsubsection*{Prínosy}

Táto práca poskytuje empirické dôkazy o~tom, ktoré prístupy k~lokálnym \gls{llm} sú prakticky použiteľné v~prostrediach terminálu, identifikuje chybové režimy klasifikácie príkazov a~agentového riadenia a~ponúka funkčného open-source asistenta založeného na \gls{rag} ako základ pre ďalší výskum. Výsledný systém je užitočný ako offline asistent pre používateľov \gls{cli}, pretože dokáže odpovedať na otázky o~lokálnych dokumentoch bez závislosti od siete. Hlavným poznatkom je, že lokálne systémy \gls{llm} sa stávajú efektívnymi vtedy, keď sú vhodne kombinované s~optimalizovanými stratégiami vyhľadávania na~úrovni pasáží.

% This thesis investigates building a~practical, offline Linux \gls{cli} assistant
% (shAI) using small local \glspl{llm} on consumer-grade hardware. It focuses on
% privacy, offline availability, and interactive latency under constrained
% resources.
%
% The thesis sets three primary goals:
% \begin{enumerate}
%   \item Analyze the landscape of existing terminal assistants and identify their
%   limitations for local use.
%   \item Design and implement a~prototype system that operates entirely on local models.
%   \item Evaluate the feasibility of key components using objective metrics.
% \end{enumerate}
%
% The resulting prototype, shAI, is open source and serves as a~reference for
% future research. The resume below condenses methods, implementation, and
% empirical results.
%
% \subsection*{Theoretical Background}
%
% This section describes the theory background relevant for the selected topic.
%
% \subsubsection*{Terminal-based Interface Concepts}
%
% This section defines interfaces relevant to our solution: \glspl{cli} and
% \glspl{tui}.
%
% \paragraph{Interface} is a~shared boundary where two separate components
% exchange information in a~certain defined way~\cite{hookway2014}.
%
% \paragraph{Hardware terminal} is an electronic device for entering and
% transcribing data, providing text-based interaction with computer
% systems~\cite{terminfo}.
%
% \paragraph{Terminal emulator} (later \textbf{Terminal}) replicates hardware
% terminal functionality as software.
%
% \paragraph{Terminal application} runs within terminal emulators and provides
% text-based user interfaces via \glspl{cli} or \glspl{tui}. Communication occurs
% through standard input/output streams, restricting interaction to
% text~\cite{man_man,linfo_terminal}.
%
% \paragraph{\gls{cli}} extends interface definitions to command-line
% interfaces --- text-based interfaces where users input commands and receive text
% responses.
%
% %TODO: fix citations for cli sec in orig
% \paragraph{Shell} is a~program providing direct system access that implements
% \gls{cli}, \gls{tui}, or \gls{gui}. \gls{cli} shells allow users to execute
% commands and \gls{cli} utilities~\cite{linfo_terminal}.
%
% \paragraph{\gls{cli} utility} is a~program executed within \gls{cli}
% environments
%
% \paragraph{\gls{tui}} extends \glspl{cli} providing structured, intuitive
% interfaces using libraries like ncurses and terminal
% capabilities~\cite{linfo_terminal}.
%
% \subsection*{\gls{llm} Concepts}
%
% %~This section creates a~background for understanding the \glspl{llm}.
% %
% %~\subsubsection*{\glspl{llm}}
%
% \paragraph{Language models} predict and generate text sequences, crucial for
% translation, summarization, and content creation, evolving with neural network
% advances~\cite{naveed2023}.
%
% \paragraph{\glspl{llm}} are language models based on transformer architecture
% with large parameter counts, trained on massive text datasets (e.g.,
% \gls{gpt}~\cite{radford2018improving}, \gls{bert}~\cite{devlin2018bert},
% \gls{lstm}~\cite{hochreiter1997lstm} or Mamba~\cite{gu2023mamba}). They
% demonstrate strong natural language understanding and complex task-solving via
% text generation.
%
% \subsubsection*{\glspl{llm} Architecture}
%
% The architecture of \glspl{llm} can be described in layers, each processing the
% input information and passing it to the next
% one~\cite{mikolov2013distributed,pennington2014glove,zhao2023, attention}.:
%
% \begin{itemize}
%     \item \textbf{Tokenizer} breaks input text into tokens, with the fixed
%     vocabulary defined at training.
%     \textsl{New or rare words may not process accurately, depending on the
%     size of the tokenizer}
%
%     \item \textbf{Embedding Layer} converts tokenized input into dense vector
%     representations (like fixed vectors or position-dependent vectors).
%
%     \item \textbf{Transformer Layers} consist of multiple
%     attention~\cite{attention} heads and feedforward networks that allow the
%     model to capture complex relationships between tokens.
%
%     \item \textbf{Output Layer} generates final output (text, probabilities) via a
%     linear layer with softmax activation.
% \end{itemize}
%
% In the original article, there are two mechanisms, built on the top of attention
% mechanism~\cite{zhao2023}: %TODO: rephrase
% \begin{itemize}
%   \item \textbf{Encoder} maps input token sequences to continuous representations
%     via multi-head self-attention and feedforward networks.
%   \item \textbf{Decoder} generates output token sequences autoregressively. It
%     includes encoder-style sub-layers plus multi-head attention over encoder
%     output.
% \end{itemize}
%
% Modern \glspl{llm} primarily use \textbf{Decoder-only} architecture, typically
% easier to scale, train and fine-tune than Encoder-Decoder models, achieving better
% performance on the same training
% tokens~\cite{zhao2023,wei2023chainofthoughtpromptingelicitsreasoning}.
%
% \paragraph{Context Window} is the maximum token count a
% Transformer processes at once. Smaller prompts process faster; larger prompts
% lose information if they exceed the window.
%
% \paragraph{\gls{moe}}
% architecture splits the feedforward network into multiple specialized expert
% subnetworks with a~router selecting only a~subset (typically one or two) for
% processing.
% \gls{moe} models are more computationally efficient but perform worse than dense
% models with the same parameter count. However, when VRAM is unconstrained,
% \gls{moe} can exceed dense model performance within equivalent computational
% budgets~\cite{dai2024deepseekmoe}.
%
%
% \subsubsection*{Embedding \gls{llm}}
%
%
% \paragraph{Embedding \gls{llm}} is a~neural model based on transformer
% architecture for generating vector representations of the documents.
% They capture more complex token relationships than traditional
% models, using attention mechanisms to weigh token importance, providing more
% accurate semantic
% representations~\cite{zhang2024systematicsurveytextsummarization,khan2026retrieval}.
%
% \subsection*{Advanced \gls{llm} Concepts}
%
% This section mentions important \gls{llm} concepts, which are commonly used in
% various \gls{llm} applications.
%
% \subsubsection*{\gls{llm} AI Agent Systems}
%
% \paragraph{AI agents} are autonomous entities capable of planning,
% decision-making, and performing actions to achieve complex
% goals~\cite{naveed2023,zhao2023}.
%
% AI Agents improve upon traditional \glspl{llm}:
% \begin{itemize}
%   \item \textbf{Hierarchical Planning:} Break complex tasks into manageable sub-tasks.
%   \item \textbf{Tool Use:} Leverage external tools, RAG pipelines, or APIs beyond text generation.
%   \item \textbf{Iterative Generation:} Generate multiple reasoning steps
%   iteratively, refining outputs and adapting to new information.
% \end{itemize}
% The workflow of typical \gls{llm} agent workflow is described by Algorithm~\ref{alg:agent}.
%
% AI Agents decompose complex objectives into manageable steps, iteratively
% prompting themselves using various tools~\cite{nvidia-llm-agent}.
%
% \subsubsection*{Retrieval-Augmented Generation (RAG)}
%
% \textbf{\gls{rag}} is a~technique that allows \glspl{llm} to utilize external
% knowledge bases, significantly reducing hallucinations and increasing
% accuracy~\cite{zhao2023}. The workflow of \gls{rag} is described by Algorithm~\ref{alg:rag}.
%
% \paragraph{Retrieval Granularities}
%
% \gls{rag} systems can retrieve information at multiple levels of granularity,
% including \textbf{document-level, passage-level, and sentence-level} retrieval.
% The choice of granularity affects the computational load for creating and
% performing retrieval operations, the accuracy of retrieval, and determines
% generated output size.
%
% The section also mentions different \gls{rag} evaluation metrics, like
% precision@k or MRR.
% %~\paragraph{RAG Evaluation Metrics}
% %
% %~Several metrics are essential for evaluating \gls{rag} system performance.
% %
% %~\textbf{Precision@k} measures the proportion of relevant documents appearing
% %~within the top-k retrieved results, defined as:
% %~\[
% %~  \text{Precision@k} =~\frac{\text{number of relevant documents in top-k}}{k}
% %~\]
% %~Higher precision at low k~values (e.g., k=1) is generally preferred, as it
% %~indicates that the system's top results are highly
% %~relevant~\cite{rag-evaluation,langsmith2026rag}.
% %
% %~\textbf{Recall@k}
% %~measures the proportion of all relevant documents successfully retrieved within
% %~the top-k results:
% %~\[
% %~  \text{Recall@k} =~\frac{\text{number of relevant documents retrieved in top-k}}{\text{total number of relevant documents}}
% %~\]
% %~Recall provides a~complementary perspective to precision, indicating whether the
% %~system finds a~significant portion of all available relevant
% %~documents~\cite{rag-evaluation,langsmith2026rag}.
% %
% %~\textbf{Mean Reciprocal Rank} measures the average rank of the first relevant
% %~result across all queries:
% %~\[
% %~  \text{MRR} =~\frac{1}{|Q|} \sum_{i=1}^{|Q|} \frac{1}{\text{rank}_i}
% %~\]
% %~where rank\textsubscript{i} is the position of the first relevant document for
% %~query i, and |Q| is the total number of queries. MRR rewards systems that find
% %~relevant documents early in the ranking and is particularly useful when users
% %~typically examine only the top few
% %~results~\cite{rag-evaluation,langsmith2026rag}.
%
% \subsubsection*{Fine-tuning}
%
% Fine-tuning is the process of taking a~pre-trained \gls{llm} and adapting it to
% a~specific task by training it on a~smaller, task-specific dataset. This
% approach allows the model to learn task-specific patterns and improve its
% performance on the target task while maintaining low training effort and
% costs~\cite{zhao2023}. While not employed in this thesis, fine-tuning is an
% important technique in \gls{llm} specialization and represents a~potential
% avenue for future enhancement.
%
% \subsection*{Document Representation And Retrieval}
%
% This section provides the theoretical background about document representation,
% storage and retrieval.
%
% \subsubsection*{Document Representation}
%
% \paragraph{Semantic representation of documents} is the process of converting
% documents into a~format that captures their conceptual and semantic meaning in a
% form suitable for machine processing. The semantic representation of documents
% is generally categorized into two primary groups~\cite{khan2026retrieval}:
% \begin{itemize}
%     \item \textbf{Sparse vector representations} employ traditional keyword-based
%     approaches such as bag-of-words (\gls{bow}), one-hot encoding, latent semantic
%     indexing (\gls{lsindexing}), or term frequency-inverse document frequency
%     (\gls{tfidf}). These are computationally efficient
%     and well-suited for simple queries, but may struggle with complex
%     queries requiring understanding of context~\cite{khan2026retrieval}.
%
%     \item \textbf{Dense vector representations} employ embedding \glspl{llm} to generate
%     embeddings that capture the semantic meaning of text. It is more
%     effective for complex queries, but requires more computational resources
%     than sparse approaches~\cite{khan2026retrieval}.
%
%     \item Additionally, \textbf{hybrid retrieval techniques} attempt to
%     combine the good sides of both by using techniques, such as the SPLADE
%     model, which enriches sparse representations with semantically related
%     terms~\cite{khan2026retrieval}.
% \end{itemize}
%
% For this research, dense retrieval is employed as the primary approach, as it is
% better suited for complex queries and can capture the semantic meaning of
% documents, which is more relevant for the task of selecting appropriate
% documents based on user queries.
%
% %~\paragraph{Dense Representation Vector Generation}
% %~Since the embedding layer represents documents as collections of vectors (one
% %~per token), techniques are needed to aggregate these into one
% %~vector~\cite{behrendt2025maxpoolbertenhancingbertclassification,devlin2018bert}.
% %
% %~\begin{itemize}
% %~    \item{\textbf{Classical aggregation approaches}} include averaging, max pooling, and
% %~    other simple methods to combine token embeddings into a~single vector. 
% %~    \item{\textbf{Neural network-based approaches}} employ trainable models such as
% %~    transformers or recurrent neural networks to generate document vectors.
% %~\end{itemize}
%
% %~\subsubsection*{Document Storage and Retrieval}
% %
% %~Once document representations are generated, they must be stored in a~manner
% %~enabling efficient similarity-based retrieval. For dense vector representations,
% %~\textbf{vector databases} are the appropriate storage mechanism.
%
% \paragraph{Vector databases} are specialized database
% systems designed to store and manage high-dimensional vector data and perform
% efficient similarity searches based on various distance metrics, including
% cosine similarity, L2 distance, and Euclidean distance. Examples of widely-used
% vector databases include FAISS (Facebook AI Similarity Search)~\cite{faiss2026},
% Annoy~\cite{annoy2026}, and Milvus~\cite{milvus2026}. These are commonly used
% for storing and retrieving document vector representations.
%
% \subsection*{Linux Manual Pages}
%
% \textbf{Linux man pages} (or \textbf{manual pages}) are processed as data for
% validation. \texttt{man} is the system's manual pager displaying program,
% utility, or function documentation~\cite{man_man}.
%
% Manual pages divide into conventional groups and sections. Key sections are
% \textbf{DESCRIPTION} (detailed program description, functionality, usage) and
% \textbf{OPTIONS} (flags, arguments, descriptions)~\cite{man_man}.
%
% \subsection*{Design and Implementation} %~Left the same
% The design requires offline operation and minimal external
% dependencies. Linux man pages are chosen as the primary knowledge source.
%
% The shAI system is implemented as a~modular application with three components:
% \begin{enumerate}
%     \item The command-vs.-query classifier (later deprioritized)
%     \item The document retriever that builds and searches embedding indices, based
%     on \gls{faiss} and Ollama for local model inference
%     \item The summarizer that produces concise responses from retrieved content
% \end{enumerate}
% Design decisions are supported by architecture and use-case diagrams (see
% Figures~\ref{fix:use-cases} and \ref{fig:component-arch}).
%
% Beyond the core components, the implementation includes user-facing workflows
% for system initialization, index building, query execution, and configuration
% changes.
%
% The classifier component is designed as a~prototype. It can be invoked on
% the current user input buffer and return a~binary decision, but integrating this
% into a~general-purpose shell has proven to be impractical, but not without
% potential, as some of existing solutions (for example, llm-zsh-plugin) have
% shown that there are application-specific ways to achieve this.
%
% The retriever component handles document gathering, embedding generation, and
% vector index storage, and it supports three search approaches:
% \begin{enumerate}
%   \item Document-level retrieval
%   \item Document retrieval with passage-level re-ranking in runtime
%   \item Passage-level retrieval
% \end{enumerate}
% It also explicit configuration for various retrieval parameters like top-k
% results, passage size, and index storage location, enabling controlled
% experiments. Indexing is performed on preprocessed man pages stored as text files
% and stored in a~local database directory. 
%
% The summarizer component provides the final user-facing response. It receives a
% query and a~set of retrieved passages, formats a~prompt with explicit structure,
% and generates a~concise answer focused on the user question. Initially intended
% to be fully functional agentic system, during the early stages of manual testing
% it has proven to be not reliable and precise, as the \gls{llm} usually failed in
% decision making, and main focus has shifted to the retriever component.
%
% The technology stack centers on Python with the OpenAI-compatible client library
% for model access, \gls{faiss} for similarity search, and standard system
% libraries for filesystem and subprocess management. Ollama provides local model
% runtime with configurable context length and GPU support. The test suite is
% config-driven, enabling reproducible runs across model variants and retrieval
% strategies.
%
% Implementation also includes a~\gls{cli} entry point that exposes configuration loading
% and test execution. Configuration files specify which model to use for embedding
% and generation, and the system is designed to reuse the same interfaces for both
% local and remote backends if needed. This modular structure allows the retrieval
% and summarization components to be tested independently.
%
% The model families used in experiments are summarized in
% Table~\ref{tab:llm-summary} and~\ref{tab:embedding-summary}).
%
%
% \subsection*{Evaluation}
%
% This section describes the evaluation process for the solutions.
%
% \subsubsection*{Natural Language vs. Shell Command Classifier}
%
% For testing the classifier, the manually curated collection\\
% \texttt{classifier-100} contains 100 inputs with a~balanced mix of natural
% language queries and shell command inputs. It is split as follows:
% \begin{itemize}
%   \item 50 natural language queries (e.g., ``How do I~list files in a
%     directory?'')
%   \item 50 shell command inputs (e.g., ``ls -la /home/user'')
% \end{itemize}
%
% \paragraph{Testing Results}{}
% Table~\ref{tab:eval-classifier-results} summarizes the classification accuracy
% and average per-query latency for each model.
%
% \paragraph{Key Findings and Limitations}
% Evaluation and testing of this component revealed that the approach for
% classification is unreliable, reaching only 50\%--87\%, with hith latencies, up
% to 2.9s, which is not reliable and fast enough for a~shell wrapper.
% Additionally, the generic shell wrapper idea itself becomes a~bottleneck
% and maintenance nightmare.
% This component remains a~\textbf{proof-of-concept} and was not
% pursued further.
%
% \subsubsection*{Document Retriever}
%
% \paragraph{Evaluation Approach}
% The retriever is evaluated on its ability to return relevant documents within
% top-k results using Precision@k and MRR, with three automated test variants:
% \begin{enumerate}
%   \item Document-level retrieval (test1)
%   \item Document retrieval with passage-level re-ranking (test2)
%       \textit{(retrieving documents first, then re-ranking passages within them, evaluating by passage matches)}
%   \item Passage-scope retrieval (test3)
% \end{enumerate}
% The test is concluded on 
% \texttt{sample-10} (14 documents, 250 queries) and
% \texttt{sample-100} (101 documents, 300 queries) test collections.
%
% \paragraph{Testing Results}{}
% Tables~\ref{tab:eval-retriever-sample10-doc}--\ref{tab:eval-retriever-sample100-passage}
% summarize document-level retrieval (test1), passage-level re-ranking (test2), and
% passage-level retrieval (test3). All tests used top\_k=5; test2 uses
% top\_k\_extended=10.
%
% \paragraph{Key Findings}
% The document retriever component proved robust and practical with all approaches
% in terms of accuracy. On \texttt{sample-10}, document-level Top-1 accuracy
% ranged from 83.6--86.8\% with Top-5 accuracy up to 100\%. On
% \texttt{sample-100}, Top-1 ranged from 62.7\% to 72.7\% with Top-5 up to 92.0\%.
% The qwen3-embedding:0.6b model provided the best precision, but its index build
% time was substantially higher on the larger corpus (128--134s vs.\ 1.6--2.3s for
% Granite embeddings). Though, it is important to note that secound and third
% approaches have much smaller output size (avg. 5K characters, which is about 4K
% tokens\ vs. avg 80K chars in first approach), which is critical for the
% summarization component, as it significantly reduces the latency and increases
% the accuracy of the generated answer.
%
% \subsubsection*{End-to-end testing}
%
% \paragraph{Evaluation Approach}
% The complete \gls{rag} system (document retrieval +~information
% summarization) is evaluated end-to-end to assess real-world
% performance including document indexing time, query processing, and
% user response latency under different hardware configurations and document
% corpus sizes. The evaluation focuses on practical usability for \gls{cli}
% assistance.
%
% \paragraph{Testing Collection}{}
% For end-to-end testing, we use query set (10 queries) and run it against the \texttt{sample-10}, \texttt{sample-100}, and \texttt{sample-1000} collections.
%
% \paragraph{Testing Results}{}
% Table~\ref{tab:eval-21-results}--\ref{tab:eval-32-results} presents sample results from evaluation
% for each collection. Values are average per-query latencies in seconds. The
% evaluation was conducted on a~system with an NVIDIA RTX 4060 8GB GPU. The testing
% is conducted for the lightest (qwen2.5:0.5b for generation and
% ibm/granite-embedding:30m for retrieval) and the heaviest (qwen3:1.7b for
% generation and qwen3-embedding:0.6b for retrieval) models, to provide a~range of
% expected performance across different hardware configurations.
%
% \paragraph{Key Findings}
% Document-level retrieval and document-level retrieval with passage reranking
% show the best performance when using small embedding models like
% ibm/granite-embedding:30m. The latency of runtime passage-level reranking is low
% enough to perform better, than passage-level approach, but keeping the output
% size small and accuracy the same. In contrast, when using bigger embedding
% models, like qwen3-embedding:0.6b, the passage-level approach outperforms the
% other approaches, as the computation time needed for vector creation for
% the whole document or runtime reranking becomes too expensive.
%
% \subsubsection*{Regression tests}
% The software regression tests that were implemented to
% ensure the stability and reliability of the ShAI system during development and
% future maintenance. The tests cover key paths of the application, including
% document retrieval, summarization, and end-to-end query processing.
% Along with the software regression tests, we also implemented a~test regression
% suite, which ensures that with any change of the system, none of the tests are
% broken (especially statistical tests, which usually need a~lot of time to run).
% This allows us to maintain the stability of the system with ease, and to quickly
% identify any issues that may arise during development.
%
% \subsubsection*{Discussion of Future Directions}
%
% Future work within the application propose such improvements:
% \begin{itemize}
%     \item \textbf{Broader task support:} test other bounded tasks such as error
%     explanation or template-based generation. 
%     \item \textbf{User interface improvements:} add richer interaction while
%     keeping latency low. 
%     \item \textbf{Software optimization:} improve caching, model loading, and
%     retrieval tuning. 
%     \item \textbf{Methodology:} compare evaluation methods and explore
%     distillation, quantization, or newer local architectures. 
% \end{itemize}
% Future work not limited to the scope of the developed application can continue
% in such directions:
% \begin{itemize}
%     \item \textbf{Comparing Different Model Testing Approaches and Frameworks}
%     \item \textbf{Model Optimization and Architectures Exploration}
%     \item \textbf{Exploring Different Harnessing Approaches}
% \end{itemize}
%
% \subsubsection*{Contributions}
%
% This work provides empirical evidence about which local \gls{llm} approaches are
% viable in terminal environments, identifies failure modes of command
% classification and agentic control, and offers a~working open-source
% \gls{rag}-based assistant as a~basis for further research.
% The resulting system is useful as an offline assistant for \gls{cli} users
% because it can answer questions about local documents without network
% dependency. The main lesson is that local \gls{llm} systems become effective
